% !TeX root = ../navier-stokes-manuscript.tex
% ============================================================
% 2.1 Vortex Stretching
% ============================================================

The momentum equation has four pieces:
\[
  \underbrace{\partial_t u}_{\text{local acceleration}}
  +
  \underbrace{(u\cdot\nabla)u}_{\text{nonlinear transport}}
  =
  \underbrace{-\nabla p}_{\text{pressure}}
  +
  \underbrace{\nu\Delta u}_{\text{viscous diffusion}}.
\]
The central regularity problem lies in the competition between nonlinear transport and viscosity.
The hope is simple: viscosity smooths the fluid faster than nonlinear transport can sharpen it.
If that remains true at every scale, the velocity stays smooth.

\subsection{Vortex Stretching}\label{sub:ThePerfectlyStretchingVortex}

Take a small material fluid element carrying vorticity along one axis. Axial strain lengthens this vortex segment; incompressibility contracts its cross-section and draws its rotating fluid inward, while the aligned stretching contribution speeds the interior rotation. Viscosity acts at the same time, so the rotation grows only when the complete vorticity balance stays positive.

Suppose the strain is locally uniform across the element and extends it along a unit direction \(e\). If \(L(t)\) is its length, write
\[
  S:=\frac12\bigl(\nabla u+(\nabla u)^{\mathsf T}\bigr),
  \qquad
  Se=\sigma(t)e,
  \qquad
  \sigma(t)>0,
  \qquad
  \frac{\dot L(t)}{L(t)}
  =
  e\cdot Se
  =
  \sigma(t).
\]
The positive strain makes \(L(t)\) increase. The element cannot lengthen without becoming narrower: incompressibility keeps its material volume \(\mathcal V\) fixed. Define its effective transverse area by \(A(t):=\mathcal V/L(t)\) and its effective radius by \(r_{\mathrm{eff}}(t):=\sqrt{A(t)/\pi}\). Then
\[
  \nabla\cdot u=0,
  \qquad
  \frac{d}{dt}(A L)=0,
  \qquad
  \frac{\dot A}{A}=-\sigma(t),
  \qquad
  \frac{\dot r_{\mathrm{eff}}}{r_{\mathrm{eff}}}
  =
  -\frac{\sigma(t)}{2}.
\]
As the segment lengthens, its material radius contracts at half the logarithmic rate.

With \(\omega:=\nabla\times u\), the vorticity carried by this segment satisfies
\[
  \frac{D\omega}{Dt}
  =
  S\omega+\nu\Delta\omega,
  \qquad
  \frac{D}{Dt}
  :=
  \partial_t+u\cdot\nabla,
\]
and under the conditional alignment \(\omega=|\omega|e\),
\[
  S\omega
  =
  \sigma(t)\omega,
  \qquad
  \frac12\frac{D|\omega|^2}{Dt}
  =
  \sigma(t)|\omega|^2
  +
  \nu\,\omega\cdot\Delta\omega.
\]
The first term records the spin-up caused by aligned stretching. The pointwise viscous contribution has no fixed sign, and net growth occurs only on intervals where the full right-hand side is positive.

Finite kinetic energy does not stop this local strengthening. The energy identity proved in \Cref{prop:ClassicalKineticEnergyIdentity} and the antisymmetric part of \(\nabla u\) give
\[
  \norm{u(t)}_2
  \leq
  \norm{u_0}_2
  <\infty,
  \qquad
  \left|\frac12\bigl(\nabla u-(\nabla u)^{\mathsf T}\bigr)\right|_{\mathrm F}
  =
  \frac{|\omega|}{\sqrt2}
  \leq
  |\nabla u|_{\mathrm F}.
\]
The velocity can retain finite energy while the rotational part of its gradient rises inside a narrowing region. The shrinking width, rather than the total kinetic energy, forces the next test: whether viscosity becomes relatively stronger when the equation is viewed at a finer scale.

\divider{\NavierStokes{} Scaling}

At a time \(t_0<T_*\), call the material radius \(\ell:=r_{\mathrm{eff}}(t_0)\). For \(\lambda>1\), the exact \NavierStokes{} comparison at width \(\lambda^{-1}\ell\) is
\[
  u_\lambda(x,t)
  :=
  \lambda u(\lambda x,\lambda^2t),
  \qquad
  p_\lambda(x,t)
  :=
  \lambda^2p(\lambda x,\lambda^2t).
\]
At time \(\lambda^{-2}t_0\), the corresponding segment in \(u_\lambda\) has effective radius \(\lambda^{-1}\ell\). This is another solution at another scale, not a later state of the original segment. Its four momentum terms transform together:
\[
\begin{aligned}
  \partial_tu_\lambda(x,t)
  &=
  \lambda^3(\partial_tu)(\lambda x,\lambda^2t),
  \\
  (u_\lambda\cdot\nabla)u_\lambda(x,t)
  &=
  \lambda^3\bigl((u\cdot\nabla)u\bigr)(\lambda x,\lambda^2t),
  \\
  \nabla p_\lambda(x,t)
  &=
  \lambda^3(\nabla p)(\lambda x,\lambda^2t),
  \\
  \nu\Delta u_\lambda(x,t)
  &=
  \lambda^3\nu(\Delta u)(\lambda x,\lambda^2t).
\end{aligned}
\]
Nonlinear transport and viscosity both gain the factor \(\lambda^3\); pressure and local acceleration gain it as well. Moving to a finer width gives viscosity no relative advantage.

\divider{Critical Height}

Although the equation preserves its balance, its Sobolev norms move in different directions. Let \(\Lambda:=(-\Delta)^{1/2}\). Since
\[
  \widehat{u_\lambda}(\xi,t)
  =
  \lambda^{-2}\widehat u(\lambda^{-1}\xi,\lambda^2t),
\]
a change of variables gives
\[
  \norm{u_\lambda(t)}_{\dot H^s}^2
  =
  \lambda^{2s-1}
  \norm{u(\lambda^2t)}_{\dot H^s}^2.
\]
At \(s=0\), kinetic energy loses one power of \(\lambda\). At \(s=1\), the first derivatives gain one. Between them the exponent vanishes at \(s=\tfrac12\), which selects
\[
  H(t)
  :=
  \frac12\norm{u(t)}_{\dot H^{1/2}}^2
  =
  \frac12\norm{\Lambda^{1/2}u(t)}_2^2.
\]
Writing \(H_\lambda\) for this quantity evaluated on the comparison field,
\[
  \norm{u_\lambda(t)}_2^2
  =
  \lambda^{-1}\norm{u(\lambda^2t)}_2^2,
  \qquad
  H_\lambda(t)=H(\lambda^2t),
  \qquad
  \norm{\nabla u_\lambda(t)}_2^2
  =
  \lambda\norm{\nabla u(\lambda^2t)}_2^2.
\]
Energy falls and the first-derivative norm rises under the comparison, while the half-derivative height remains unchanged. This critical height is a global functional of the velocity field, not a height attached to the material vortex. Scale covariance selects it, but does not determine whether it rises or falls along the original solution; that depends on the sign of its time derivative.

\divider{Nonlinear Production versus Viscous Dissipation}

At critical height, the nonlinear term contributes the signed production
\[
  \mathcal P_{1/2}[u](t)
  :=
  -\operatorname{Re}
  \pair
    {\Lambda^{1/2}u(t)}
    {\Lambda^{1/2}\bigl((u(t)\cdot\nabla)u(t)\bigr)}_{L^2}.
\]
The pressure pairing vanishes because \(\nabla\cdot u=0\), and the Laplacian contributes \(-\nu\norm{\Lambda^{3/2}u}_2^2\). Thus
\[
  H'(t)
  +
  \nu\norm{\Lambda^{3/2}u(t)}_2^2
  =
  \mathcal P_{1/2}[u](t).
\]
The critical height rises exactly when signed nonlinear production exceeds simultaneous viscous dissipation. Under the same rescaling,
\[
  \mathcal P_{1/2}[u_\lambda](t)
  =
  \lambda^2\mathcal P_{1/2}[u](\lambda^2t),
  \qquad
  \nu\norm{\Lambda^{3/2}u_\lambda(t)}_2^2
  =
  \lambda^2\nu\norm{\Lambda^{3/2}u(\lambda^2t)}_2^2.
\]
Both rates gain the same square factor. Their relative strength remains unresolved at finer scale, and this covariance does not tell us how long the original vortex spends at any width. A width-time relation has to come from a separate calculation.

\divider{The Heat Clock}

For the viscous term alone, the heat equation \(v_t=\nu\Delta v\) evolves each Fourier mode by
\[
  \widehat v(\xi,t+\delta t)
  =
  e^{-\nu|\xi|^2\delta t}\widehat v(\xi,t).
\]
If vorticity is localized across a transverse width \(r\) and has active frequencies \(|\xi|\simeq r^{-1}\), its heat evolution changes by order one when
\[
  \nu|\xi|^2\delta t\simeq1,
\]
which yields the linear comparison time
\[
  \tau_\nu(r)
  :=
  \frac{r^2}{\nu}.
\]
This clock belongs to scale-localized heat modes. The material radius \(r_{\mathrm{eff}}\) alone does not establish vorticity localization, active frequency, or the lifetime of the vortex. Under the stated localization hypothesis, halving the width quarters the clock; more generally,
\[
  \tau_\nu(\lambda^{-1}r)
  =
  \lambda^{-2}\tau_\nu(r).
\]
Successively finer localized widths carry successively shorter comparison times.

\divider{The Zeno Cascade}

Repeat the arithmetic halving. For
\[
  r_n:=2^{-n}r_0,
  \qquad
  \tau_n:=\frac{r_n^2}{\nu},
\]
the heat clocks satisfy
\[
  \tau_n
  =
  \frac{r_0^2}{\nu}4^{-n},
  \qquad
  \sum_{n=0}^{\infty}\tau_n
  =
  \frac{4r_0^2}{3\nu}
  <
  \infty.
\]
The finite sum is arithmetic, not a realized cascade. To turn it into a candidate singular history, one \NavierStokes{} solution must carry the original material segment through widths
\[
  r_0>r_1>r_2>\cdots\longrightarrow0
\]
at times
\[
  t_0<t_1<t_2<\cdots\uparrow T_*<\infty,
  \qquad
  t_{n+1}-t_n\asymp\tau_n,
\]
with constants independent of \(n\). More precisely, at each \(t_n\), the segment radius \(r_{\mathrm{eff}}(t_n)\), the transverse width across which vorticity remains localized around that segment, and \(r_n\) must be mutually comparable with uniform constants. The active vorticity frequencies must also be comparable to \(r_n^{-1}\), and the time gap must remain comparable to \(r_n^2/\nu\).

Under those same-history hypotheses, the remaining comparison time satisfies
\[
  T_*-t_n
  \asymp
  \sum_{k=n}^{\infty}\tau_k
  =
  \frac{4r_n^2}{3\nu}.
\]
Both the next transition and the time left then shrink like \(r_n^2/\nu\). The arithmetic can describe a candidate history only if one velocity field produces every successive localized narrowing of the tracked vortex on those collapsing comparison intervals.
